Supersymmetry~(SUSY)~\cite{Golfand:1971iw,Volkov:1973ix,Wess:1974tw,Wess:1974jb,Ferrara:1974pu,Salam:1974ig} is a theoretical framework based on a space-time symmetry that predicts the existence of 
a partner for each Standard Model~(SM) particle. These partner particles have the same quantum numbers as their SM counterparts, but differ by one half unit of spin. SUSY offers a solution to the hierarchy problem~\cite{Sakai:1981gr,Dimopoulos:1981yj,Ibanez:1981yh,Dimopoulos:1981zb} if some of the partner particles have masses around the electroweak scale. 
The supersymmetric partners are named to reflect their SM counterparts: for example, the $\tau$-slepton is the SUSY partner of the $\tau$-lepton. 
The supersymmetric partners of the electroweak gauge bosons and the Higgs bosons, collectively referred to as electroweakinos, consist of the bino and winos and the higgsinos, respectively.
These states mix to form the mass eigenstates: 
the bino, neutral wino, and neutral higgsinos mix to form the neutralinos ($\tilde{\chi}_{1,2,3,4}^{0}$), while the charged wino and charged higgsinos mix to form charginos ($\tilde{\chi}_{1,2}^{\pm}$), where the subscripts denote increasing mass.
The superpartners of the third-generation charged leptons ($\tau_{\mathrm{L}}, \tau_{\mathrm{R}}$) are referred to as $\tau$-sleptons ($\stau_{\mathrm{L}}, \stau_{\mathrm{R}}$) and mix to form two mass eigenstates $\stau_{1}$ and $\stau_{2}$, where $\stau_{1}$ is the lighter of the two and is henceforth referred to as the $\tau$-slepton ($\stau$).
 $R$-parity-conserving supersymmetry~\cite{Farrar:1978xj} addresses several limitations of the SM, providing a natural candidate for dark matter~(DM) in the form of the lightest supersymmetric particle~(LSP)~\cite{Goldberg:1983nd,Ellis:1983ew}. 
 In SUSY scenarios with very heavy, strongly interacting superpartners, the direct production cross-section of electroweakinos and sleptons in high-energy proton-proton (\pp) collisions at the Large Hadron Collider~(LHC) can significantly exceed that of strongly interacting superpartners. 

In scenarios where the \ninoone and \chinoonepm are wino-like and all other SUSY particles are decoupled, the mass difference, $\mdiffchi$, is small ($\sim$100\,\MeV) and arises from SM radiative corrections~\cite{Ibe:2012sx,Fukuda:2017jmk}.
Anomaly-mediated supersymmetry-breaking~(AMSB) models~\cite{Giudice:1998xp,Randall:1998uk} naturally predict such mass splittings and give rise to a wino-like LSP.
In these scenarios, the mass-splitting between the charged wino and neutral wino is suppressed at tree level by the approximate custodial symmetry, 
and has been calculated at the two-loop level to be approximately $160\,\MeV$~\cite{Ibe:2012sx}, corresponding to a chargino lifetime 
of $\tau(\chinoonepm) \sim 0.2$~ns and a decay length of $c\tau(\chinoonepm) \sim 60$~mm. 
Furthermore, a number of \enquote{natural} SUSY models~\cite{Papucci:2011wy,Barbieri:1987fn,deCarlos:1993yy}, where the \ninotwo is no longer decoupled, predict a light higgsino LSP with a mass near the electroweak scale.
In these scenarios, the higgsino mass parameter $|\mu|$ is small compared to the other electroweakino mass parameters, which could mean
higgsinos are discoverable at the LHC. While the charged and neutral higgsino states are mass degenerate at tree level, higher-order SM loop corrections generate a mass splitting of approximately $300\,\MeV$.
For chargino masses between $91\,\GeV$ and $1000\,\GeV$, the mass splitting ranges approximately from $280\,\MeV$ to $350\,\MeV$~\cite{Thomas:1998wy}, 
corresponding to chargino lifetimes of $\tau(\chinoonepm) \sim 0.026{-}0.048$~ns ($c\tau(\chinoonepm) \sim 7{-}14$~mm)~\cite{Fukuda:2017jmk}. Due to the small mass splitting, the \chinoonepm is assumed to decay via $\chinoonepm \rightarrow \pi^\pm\ninoone$ with a branching fraction of 100\% in the wino scenarios and 95\% in the higgsino scenarios~\cite{Nagata:2014wma}. The neutralino masses are set such that $m(\ninotwo) = m(\ninoone)$, hence the \ninotwo cannot decay to the \chinoonepm or \ninoone and is treated as a stable particle that escapes the detector, similarly to the \ninoone. 

Models featuring light $\tau$-sleptons can yield a dark-matter relic density consistent with cosmological observations~\cite{Vasquez:2011} and, more generally, light sleptons can participate in the co-annihilation of neutralinos~\cite{Belanger:2004ag,King:2007vh}. Long-lived $\tau$-sleptons arise naturally in co-annihilation dark-matter models~\cite{Ellis:1998kh,Jittoh:2005pq,Kaneko:2008re} and in gauge-mediated supersymmetry-breaking~(GMSB) 
models~\cite{Dine:1981gu,AlvarezGaume:1981wy,Nappi:1982hm}.
Two simplified SUSY models are considered for $\tau$-sleptons in this paper. The first is inspired by the Constrained Minimal Supersymmetric Standard Model (CMSSM)~\cite{Chamseddine:1982jx,Barbieri:1982eh,Kane:1993td}, where a long-lived $\tau$-slepton can arise from a small (${\sim}2$\,GeV) mass-splitting between the $\tau$-slepton and the $\ninoone$ LSP, CP-violating phases, and the mixing angle between the left- and right-handed $\tau$-slepton states. The CMSSM-inspired scenarios are parameterised by the $\tau$-slepton mass and lifetime, and the $\tau$-slepton decays with 100\% branching fraction as $\stau\rightarrow\tau\ninoone$.
The second scenario is inspired by GMSB models, where a long-lived $\tau$-slepton can arise from small couplings to its decay products. The LSP is the almost massless gravitino ($\tilde{G}$), the superpartner of a proposed graviton. This model is parameterised by the $\tau$-slepton mass and lifetime, and assumes a 100\% branching fraction for $\stau\rightarrow\tau\gravino$. 

In all considered scenarios, the charged SUSY particle (wino, higgsino or $\tau$-slepton) can be produced with large momentum and live long enough to traverse multiple layers of the ATLAS pixel detector before decaying. In models containing charginos, the chargino deposits energy in the innermost tracking layers before it decays and can be reconstructed as a short track if at least three layers of the tracking detector are hit.
The weakly interacting LSP will escape detection and lead to missing transverse momentum ($\vec{p}_\text{T}^\text{miss}$) and the charged pion from the chargino decay has a very low momentum and is not reconstructed as a track by existing algorithms, this gives rise to the characteristic \enquote{disappearing track} signature where the track from the chargino stops a short distance into the detector. The algorithms used to reconstruct the short tracks, referred to as tracklets, need energy deposits (hits) in at least three pixel layers. Therefore successful reconstruction requires charginos to have transverse path lengths that at least exceed the radius of the third pixel layer at  $\sim$88.5~mm.
This makes the higgsino-like scenarios considerably more challenging than the wino-like models from an experimental perspective because of the extremely short chargino lifetime predicted in the higgsino models. A similar scenario occurs in the case of $\tau$-slepton production, where a long-lived $\tau$-slepton particle will traverse the inner layers of the pixel detector before decaying into a $\tau$-lepton and the LSP. The LSP escapes the detector unobserved, resulting in missing transverse momentum, and the $\tau$-leptons evade standard reconstruction techniques, again resulting in a signature containing a disappearing track and missing transverse momentum. 

In addition to the SUSY-specific wino and higgsino models studied in this paper, the disappearing-track signature is typical of a large class of DM models that predict a DM thermal relic from a massive particle with only electroweak gauge interactions~\cite{Cirelli:2005uq}. 
The wino and higgsino models probed in this search are part of a more generic class of models containing spin-1/2 particles transforming under SU(2) symmetry, which give rise to nearly mass-degenerate doublets ($\chinoonepm$, $\ninoone$) and triplets ($\chinoonepm$, $\ninotwo$, $\ninoone$) of new particles, respectively. In both cases, the $\ninoone$ is a DM candidate, while the $\chinoonepm$ gives rise to the disappearing-track signature. These DM models have gained interest in the wider community as an important area of study at future colliders~\cite{CidVidal:2018eel,Saito:2019rtg,Capdevilla:2021fmj,Capdevilla:2024bwt}.

This paper targets two electroweak production processes, the production of charginos and neutralinos, and the production of $\tau$-slepton pairs in CMSSM- and GMSB-inspired scenarios, as shown in Figure~\ref{fig:SUSYFeynman}, in the context of simplified SUSY models~\cite{Alwall:2008ve,Alwall:2008ag,Alves:2011wf}.
In all scenarios, the charged SUSY partner is long-lived and can be reconstructed from hits in the ATLAS pixel detector. 
A high-momentum jet from initial-state radiation~(ISR) is required in order to ensure that significant missing transverse momentum is available for event triggering.

In order to reconstruct as many charginos as possible across a wide range of lifetimes, the shortest possible
track is desired. Since a helical trajectory in 3-dimensions needs at least three measurements to seed its
trajectory, the minimum track length that is possible extends out to the third pixel layer (88.5 mm). This
section describes the dedicated tracking configuration implemented to reconstruct tracklets with as few as
three pixel hits. An efficient track extension is developed for this configuration to allow the reconstruction
of longer tracks with additional pixel or SCT hits.

\begin{figure}[t!]
 \centering
  \begin{subfigure}[t]{0.3\textwidth}
    \includegraphics[width=.8\textwidth]{../figures/intro/C1N1.pdf}
    \caption{}
    \label{fig:intro:feynman:c1n1}
    \end{subfigure}
  \begin{subfigure}[t]{0.3\textwidth}
    \includegraphics[width=\textwidth]{../figures/intro/staustau-jtautauN1N1.pdf}
    \caption{}    
    \label{fig:intro:feynman:ststCMSSM}
    \end{subfigure}
  \begin{subfigure}[t]{0.3\textwidth}
    \includegraphics[width=\textwidth]{../figures/intro/staustau-jtautauG1G1.pdf}
    \caption{}
    \label{fig:intro:feynman:ststGMSB}
    \end{subfigure}
% \vspace{-0.75cm}
\caption{Representative signal diagrams for \protect\subref{fig:intro:feynman:c1n1} the electroweak production of $\chinoonepm\ninoone$, 
  and $\tau$-slepton pair production in \protect\subref{fig:intro:feynman:ststCMSSM} CMSSM-inspired or \protect\subref{fig:intro:feynman:ststGMSB} GMSB-inspired SUSY models.
  The signal signature consists of a long-lived \chinoonepm or $\tau$-slepton, missing transverse momentum, and quarks or gluons, which are observed as jets ($j$). When considering the electroweakino scenarios, the production of $\chinoonep\chinoonem$ is also considered, as is $\chinoonepm\ninotwo$ specifically in the higgsino scenario.}
  \label{fig:SUSYFeynman}
\end{figure}

Previous ATLAS searches for long-lived supersymmetric charged particles searched for events with a disappearing track with at least four pixel hits~\cite{SUSY-2018-19,SUSY-2016-06,ATL-PHYS-PUB-2017-019}. Their results benefited from the addition of a new innermost pixel tracking layer installed during the LHC long shutdown between Run~1 and Run~2. The extra layer of pixels enabled the reconstruction of tracks shorter than those in the previous Run-1 analysis~\cite{SUSY-2013-01} and enhanced the sensitivity to shorter charged-particle lifetimes. The Run-2 ATLAS search excluded wino-like charginos with lifetimes of $0.2\,$ns for masses up to $460\,\GeV$, and higgsino-like charginos with masses up to $152\,\GeV$~\cite{SUSY-2018-19}. Searches performed by the CMS Collaboration excluded wino-like charginos with masses below $474\,\GeV$ and a lifetime of $0.2\,$ns~\cite{CMS-EXO-19-010}. 
Both the ATLAS and CMS collaborations have performed searches for stable charginos and $\tau$-sleptons using ionisation-loss or time-of-flight techniques, excluding charginos with masses smaller than $1.3\,\TeV$ for lifetimes exceeding $100\,$ns, and $\tau$-sleptons with masses $200{-}560\,\GeV$ for lifetimes of $10\,$ns~\cite{HMBS-2024-68,CMS-EXO-18-002}.

%, a high momentum jet, and large missing transverse momentum
%generated by the weakly interacting stable \ninoone particles. These searches primarily focused scenarios where
%the favored lifetimes of the chargino is approximately a factor of four or more than the pure higgsino states.
%The results were interpreted in a pure higgsino scenario, but given the smaller higgsino production rates
%and relatively long length of track used in previous searches, sensitivity only extended up to 152 GeV [4]
%for the pure higgsino case. 
In order to improve on these results, significant developments were made in the tracking techniques necessary for shorter lifetime scenarios.
This search employs a dedicated tracking algorithm to reconstruct tracklets with as few as three pixel hits. 
Furthermore, to enhance sensitivity to the electroweakino production scenario, the charged pion from the decay of the $\chinoonepm$ is targeted with a specially optimised reconstruction algorithm.
Due to the small mass-splitting between the $\chinoonepm$ and $\ninoone$, the charged pion has extremely low momentum, ranging from $250$ to $350\,\MeV$ for the \chinoonepm lifetimes considered, and typically has large impact parameters measured relative to the interaction region.
In this paper, the sensitivity to charginos with wino and higgsino lifetimes motivated by SM loop corrections is significantly improved due 
to new analysis methods including updated signal region selection criteria, the previously mentioned reconstruction algorithms, and improved track-quality requirements.
The new track-quality criteria requires the disappearing track not to be associated with any energy deposits in the calorimeter.
This significantly enhances the rejection of dominant backgrounds
and accounts for two-thirds of the total improvement in the mass reach of the search in the electroweak wino-like production channel.


The paper is structured as follows: Section~\ref{sec:detector} introduces the ATLAS detector apparatus; Section~\ref{sec:DataAndMC} describes the dataset and the Monte Carlo simulation samples used; Section~\ref{sec:Reco} discusses the object-level event reconstruction, while Section~\ref{sec:Analysis} describes the signal selection optimisation and Section~\ref{sec:bkgs} describes the background estimation techniques; Section~\ref{sec:Systs} provides an overview of the systematic uncertainties considered, with Section~\ref{sec:result} presenting the results and interpretation thereof; and finally, Section~\ref{sec:conclusion} summarises the conclusions. 

